\chapter{Avaliação entre avaliadores}

\eyebrow{escores independentes e cegos + concordância}

\vspace{4pt}
A visão de \textbf{Avaliação} de um lote reúne graus de múltiplos avaliadores, cada
um pontuando de forma independente e cega, e mede o quanto eles concordam. Esta é
a referência humana que a pesquisa exige --- um grau atribuído \emph{sobre a imagem},
com uma medida de sua confiabilidade.

\section{Como funciona uma sessão}

\begin{enumerate}[leftmargin=1.4em,itemsep=3pt]
  \item Adicione ou escolha um \textbf{avaliador}.
  \item Inicie uma \textbf{sessão cega} no lote --- os escores dos demais avaliadores
        ficam ocultos, de modo que cada avaliação seja independente.
  \item Para cada carcaça com imagem, pontue dois eixos em uma escala de 1--5:
        \textbf{Conformação} (forma e desenvolvimento muscular, observáveis na
        imagem) e \textbf{Acabamento} (cobertura de gordura).
\end{enumerate}

\section{Concordância}

Assim que pelo menos dois avaliadores tiverem pontuado o mesmo conjunto de carcaças, o
painel do Projeto reporta a concordância entre avaliadores com o \textbf{$\kappa$ de Fleiss}
por eixo, juntamente com um grau de consenso. Se a concordância for insuficiente, o consenso
é sinalizado --- exatamente a salvaguarda de que a pesquisa precisa antes de confiar em um grau
como referência.

\begin{note}[title={Dois eixos, duas confianças diferentes}]
A \emph{Conformação} é observável na imagem. O \emph{Acabamento} depende da cobertura
de gordura, uma propriedade de profundidade que uma única projeção de superfície captura apenas
parcialmente --- a interface registra isso, e a validação preditiva completa do acabamento depende
da referência física medida nos mesmos animais.
\end{note}
